% Options for packages loaded elsewhere
\PassOptionsToPackage{unicode}{hyperref}
\PassOptionsToPackage{hyphens}{url}
\PassOptionsToPackage{dvipsnames,svgnames,x11names}{xcolor}
\documentclass[
  11pt,
]{article}
\usepackage{xcolor}
\usepackage[margin=1in]{geometry}
\usepackage{amsmath,amssymb}
\setcounter{secnumdepth}{5}
\usepackage{iftex}
\ifPDFTeX
  \usepackage[T1]{fontenc}
  \usepackage[utf8]{inputenc}
  \usepackage{textcomp} % provide euro and other symbols
\else % if luatex or xetex
  \usepackage{unicode-math} % this also loads fontspec
  \defaultfontfeatures{Scale=MatchLowercase}
  \defaultfontfeatures[\rmfamily]{Ligatures=TeX,Scale=1}
\fi
\usepackage{lmodern}
\ifPDFTeX\else
  % xetex/luatex font selection
\fi
% Use upquote if available, for straight quotes in verbatim environments
\IfFileExists{upquote.sty}{\usepackage{upquote}}{}
\IfFileExists{microtype.sty}{% use microtype if available
  \usepackage[]{microtype}
  \UseMicrotypeSet[protrusion]{basicmath} % disable protrusion for tt fonts
}{}
\makeatletter
\@ifundefined{KOMAClassName}{% if non-KOMA class
  \IfFileExists{parskip.sty}{%
    \usepackage{parskip}
  }{% else
    \setlength{\parindent}{0pt}
    \setlength{\parskip}{6pt plus 2pt minus 1pt}}
}{% if KOMA class
  \KOMAoptions{parskip=half}}
\makeatother
\usepackage{color}
\usepackage{fancyvrb}
\newcommand{\VerbBar}{|}
\newcommand{\VERB}{\Verb[commandchars=\\\{\}]}
\DefineVerbatimEnvironment{Highlighting}{Verbatim}{commandchars=\\\{\}}
% Add ',fontsize=\small' for more characters per line
\newenvironment{Shaded}{}{}
\newcommand{\AlertTok}[1]{\textcolor[rgb]{1.00,0.00,0.00}{\textbf{#1}}}
\newcommand{\AnnotationTok}[1]{\textcolor[rgb]{0.38,0.63,0.69}{\textbf{\textit{#1}}}}
\newcommand{\AttributeTok}[1]{\textcolor[rgb]{0.49,0.56,0.16}{#1}}
\newcommand{\BaseNTok}[1]{\textcolor[rgb]{0.25,0.63,0.44}{#1}}
\newcommand{\BuiltInTok}[1]{\textcolor[rgb]{0.00,0.50,0.00}{#1}}
\newcommand{\CharTok}[1]{\textcolor[rgb]{0.25,0.44,0.63}{#1}}
\newcommand{\CommentTok}[1]{\textcolor[rgb]{0.38,0.63,0.69}{\textit{#1}}}
\newcommand{\CommentVarTok}[1]{\textcolor[rgb]{0.38,0.63,0.69}{\textbf{\textit{#1}}}}
\newcommand{\ConstantTok}[1]{\textcolor[rgb]{0.53,0.00,0.00}{#1}}
\newcommand{\ControlFlowTok}[1]{\textcolor[rgb]{0.00,0.44,0.13}{\textbf{#1}}}
\newcommand{\DataTypeTok}[1]{\textcolor[rgb]{0.56,0.13,0.00}{#1}}
\newcommand{\DecValTok}[1]{\textcolor[rgb]{0.25,0.63,0.44}{#1}}
\newcommand{\DocumentationTok}[1]{\textcolor[rgb]{0.73,0.13,0.13}{\textit{#1}}}
\newcommand{\ErrorTok}[1]{\textcolor[rgb]{1.00,0.00,0.00}{\textbf{#1}}}
\newcommand{\ExtensionTok}[1]{#1}
\newcommand{\FloatTok}[1]{\textcolor[rgb]{0.25,0.63,0.44}{#1}}
\newcommand{\FunctionTok}[1]{\textcolor[rgb]{0.02,0.16,0.49}{#1}}
\newcommand{\ImportTok}[1]{\textcolor[rgb]{0.00,0.50,0.00}{\textbf{#1}}}
\newcommand{\InformationTok}[1]{\textcolor[rgb]{0.38,0.63,0.69}{\textbf{\textit{#1}}}}
\newcommand{\KeywordTok}[1]{\textcolor[rgb]{0.00,0.44,0.13}{\textbf{#1}}}
\newcommand{\NormalTok}[1]{#1}
\newcommand{\OperatorTok}[1]{\textcolor[rgb]{0.40,0.40,0.40}{#1}}
\newcommand{\OtherTok}[1]{\textcolor[rgb]{0.00,0.44,0.13}{#1}}
\newcommand{\PreprocessorTok}[1]{\textcolor[rgb]{0.74,0.48,0.00}{#1}}
\newcommand{\RegionMarkerTok}[1]{#1}
\newcommand{\SpecialCharTok}[1]{\textcolor[rgb]{0.25,0.44,0.63}{#1}}
\newcommand{\SpecialStringTok}[1]{\textcolor[rgb]{0.73,0.40,0.53}{#1}}
\newcommand{\StringTok}[1]{\textcolor[rgb]{0.25,0.44,0.63}{#1}}
\newcommand{\VariableTok}[1]{\textcolor[rgb]{0.10,0.09,0.49}{#1}}
\newcommand{\VerbatimStringTok}[1]{\textcolor[rgb]{0.25,0.44,0.63}{#1}}
\newcommand{\WarningTok}[1]{\textcolor[rgb]{0.38,0.63,0.69}{\textbf{\textit{#1}}}}
\usepackage{longtable,booktabs,array}
\newcounter{none} % for unnumbered tables
\usepackage{calc} % for calculating minipage widths
% Correct order of tables after \paragraph or \subparagraph
\usepackage{etoolbox}
\makeatletter
\patchcmd\longtable{\par}{\if@noskipsec\mbox{}\fi\par}{}{}
\makeatother
% Allow footnotes in longtable head/foot
\IfFileExists{footnotehyper.sty}{\usepackage{footnotehyper}}{\usepackage{footnote}}
\makesavenoteenv{longtable}
\setlength{\emergencystretch}{3em} % prevent overfull lines
\providecommand{\tightlist}{%
  \setlength{\itemsep}{0pt}\setlength{\parskip}{0pt}}
\usepackage{etoolbox}
\usepackage[htt]{hyphenat}
\AtBeginEnvironment{longtable}{\small\setlength{\tabcolsep}{4pt}}
\emergencystretch=3em
\setmonofont{DejaVuSansMono.ttf}[Scale=0.86]
\DeclareMathOperator*{\argmin}{arg\,min}
\newcommand{\llbracket}{[\![}
\newcommand{\rrbracket}{]\!]}
\usepackage{bookmark}
\IfFileExists{xurl.sty}{\usepackage{xurl}}{} % add URL line breaks if available
\urlstyle{same}
\hypersetup{
  pdftitle={Supplementary Material for AgentISA: Formal Instruction Set Design for LLM Agents},
  pdfauthor={Anonymous Author(s)},
  colorlinks=true,
  linkcolor={blue},
  filecolor={Maroon},
  citecolor={Blue},
  urlcolor={blue},
  pdfcreator={LaTeX via pandoc}}

\title{Supplementary Material for AgentISA: Formal Instruction Set
Design for LLM Agents}
\author{Anonymous Author(s)}
\date{}

\begin{document}
\maketitle
\begin{abstract}
This document is the supplementary material for \emph{AgentISA: Formal
Instruction Set Design for LLM Agents}. It contains the appendices
referenced by the main paper: Appendix A (the Turing-completeness
construction for Theorem 3.1), Appendix B (the type-safety proof for
Theorem 3.2), and Appendix C (the semantic-preservation lemmas and
proofs of Theorems 3.4--3.6, with the scope remarks of C.7). Section and
theorem numbers follow the main paper.
\end{abstract}

\appendix

\section{Proof of Theorem 3.1 (Turing
Completeness)}\label{proof-of-theorem-3.1-turing-completeness}

\subsection{Strategy}\label{strategy}

We prove that AgentISA can simulate any program in the WHILE language
over the natural numbers. WHILE programs compute exactly the partial
recursive functions---WHILE is Turing equivalent {[}43, 44{]}---so a
sound simulation of WHILE by AgentISA establishes Turing completeness.
The argument is a sufficiency (simulation) argument over instruction
classes: the instruction classes of Definition 3.1 must jointly provide
unbounded storage, sequential composition, conditional branching, and
unbounded iteration. We show the control-flow class \(\mathcal{C}\) and
the state-management class \(\mathcal{S}\) suffice; neither effectful
class (the LLM-interaction class \(\mathcal{L}\) nor the resource class
\(\mathcal{R}\)) is needed, making the result conservative with respect
to the full instruction set.

\subsection{Source Language}\label{source-language}

Let WHILE be the language

\[
P ::= x := x+1 \mid x := x-1 \mid x := y \mid \texttt{skip} \mid \texttt{if } x = 0 \texttt{ then } P_1 \texttt{ else } P_2 \mid \texttt{while } x \neq 0 \texttt{ do } P \mid P_1 ; P_2
\]

over variables ranging over \(\mathbb{N}\) (with predecessor cut off at
0), with the standard small-step semantics. Configurations are pairs
\(\langle P, \sigma \rangle\) with store
\(\sigma : Var \to \mathbb{N}\).

\subsection{Encoding}\label{encoding}

\textbf{(i) Store.} A WHILE store \(\sigma\) is represented by the
AgentISA state \(s_\sigma \in \Sigma\), in the component managed by the
state-management instructions. Each variable \(x\) occupies a dedicated
cell; cells hold unbounded naturals, so \(\Sigma\) provides an unbounded
store---the crucial ingredient of the sufficiency argument.

\textbf{(ii) Instruction semantics (as used here).} The state-management
component of \(\Sigma\) carries, alongside the variable cells, an
\textbf{evaluation stack}; the five encoding instructions below operate
on that stack. Five instructions suffice for the encoding, with the
following operational meaning:

\begin{itemize}
\tightlist
\item
  \texttt{load} \(x\): pushes the contents of the cell for \(x\) onto
  the evaluation stack;
\item
  \texttt{store} \(x\): pops the top of the evaluation stack into the
  cell for \(x\);
\item
  \texttt{update} \((f)\): pops the top of the stack, applies the unary
  arithmetic function \(f\), and pushes the result;
  \texttt{update}\((-1)\) denotes the \emph{truncated predecessor}
  \(\lambda n.\,\max(n-1, 0)\), matching WHILE's cutoff semantics on
  \(\mathbb{N}\);
\item
  \texttt{branch} \((b, Q_1, Q_2)\): pops the top of the stack; if
  predicate \(b\) holds, execute \(Q_1\), else execute \(Q_2\);
\item
  \texttt{loop} \((Q)\): execute \(Q\); if its final \texttt{branch}
  returns to the loop head, repeat, else exit---that is, \texttt{loop}
  repeats its body for as long as the body's trailing test succeeds.
\end{itemize}

\textbf{(iii) Statements.} Each WHILE statement compiles to the AgentISA
sequence below (notation defined after the table):

\begin{longtable}[]{@{}
  >{\raggedright\arraybackslash}p{(\linewidth - 2\tabcolsep) * \real{0.4286}}
  >{\raggedright\arraybackslash}p{(\linewidth - 2\tabcolsep) * \real{0.5714}}@{}}
\caption{Compilation of WHILE statements to AgentISA sequences
(\(\hat{P}\) is the compilation of \(P\)).}\tabularnewline
\toprule\noalign{}
\begin{minipage}[b]{\linewidth}\raggedright
WHILE construct
\end{minipage} & \begin{minipage}[b]{\linewidth}\raggedright
AgentISA compilation
\end{minipage} \\
\midrule\noalign{}
\endfirsthead
\toprule\noalign{}
\begin{minipage}[b]{\linewidth}\raggedright
WHILE construct
\end{minipage} & \begin{minipage}[b]{\linewidth}\raggedright
AgentISA compilation
\end{minipage} \\
\midrule\noalign{}
\endhead
\bottomrule\noalign{}
\endlastfoot
\(x := y\) & \texttt{load} \(y\); \texttt{store} \(x\) \\
\(x := x+1\) & \texttt{load} \(x\); \texttt{update}\((+1)\);
\texttt{store} \(x\) \\
\(x := x-1\) & \texttt{load} \(x\); \texttt{update}\((-1)\);
\texttt{store} \(x\) \\
\texttt{if} \(x = 0\) \texttt{then} \(P_1\) \texttt{else} \(P_2\) &
\texttt{load} \(x\); \texttt{branch}\((= 0?, \hat{P}_1, \hat{P}_2)\) \\
\texttt{while} \(x \neq 0\) \texttt{do} \(P\) &
\texttt{loop}(\texttt{load} \(x\);
\texttt{branch}\((\neq 0?, \hat{P}, \epsilon)\)) \\
\(P_1 ; P_2\) & \(\hat{P}_1 \circ \hat{P}_2\) (sequential
composition) \\
\end{longtable}

where \(\hat{P}\) denotes the compilation of \(P\). The no-op
\texttt{skip} compiles to the identity instruction \(\epsilon\); the
\texttt{while} row's trailing \texttt{branch} exits to \(\epsilon\) when
the guard fails, which is exactly the unfolding identity below.

\subsection{Simulation Lemma}\label{simulation-lemma}

\textbf{Lemma A.1.} For every WHILE configuration
\(\langle P, \sigma \rangle\) and its encoded AgentISA instruction
sequence \(\hat{P}\) started in state \(s_\sigma\):

\begin{enumerate}
\def\labelenumi{(\alph{enumi})}
\item
  if \(\langle P, \sigma \rangle \to \langle P', \sigma' \rangle\), then
  the AgentISA machine, started at configuration \((\hat{P}, s_\sigma)\)
  (remaining instruction sequence, \(\Sigma\)-state), reaches the
  configuration \((\hat{P}', s_{\sigma'})\)---remaining sequence
  \(\hat{P}'\) the encoding of the residual program \(P'\), state
  \(s_{\sigma'}\) the encoding of \(\sigma'\)---in finitely many steps;
  and
\item
  if \(P = \texttt{skip}\) (terminal), then \(\hat{P}\) terminates.
\end{enumerate}

\emph{Proof.} By structural induction on \(P\):

\begin{itemize}
\item
  \textbf{Assignments} (base cases). Each compiled sequence performs
  exactly the corresponding store update; (a) holds by executing the two
  or three instructions in order, (b) trivially.
\item
  \textbf{Conditional.} The \texttt{branch} instruction evaluates the
  zero test on the loaded value and transfers control to one of the
  compiled branch bodies. By the induction hypothesis applied to \(P_1\)
  and \(P_2\), each branch simulates its WHILE counterpart.
\item
  \textbf{Loop.} The WHILE unfolding identity
  \texttt{while\ x\ ≠\ 0\ do\ P\ ≡\ if\ x\ =\ 0\ then\ skip\ else\ (P\ ;\ while\ x\ ≠\ 0\ do\ P)}
  is simulated by the \texttt{loop} instruction's iteration: each
  iteration loads \(x\) and executes the trailing \texttt{branch}, which
  dispatches to \(\hat{P}\) when \(x \neq 0\) and to \(\epsilon\) (loop
  exit) when \(x = 0\). By an outer induction on the number of executed
  iterations (in addition to the structural induction on \(P\) announced
  above---the induction hypothesis applied to \(P\) covers each body
  execution, and the trailing branch discharges the loop
  condition)---the \(k\)-th WHILE iteration is matched by the \(k\)-th
  \texttt{loop} iteration.
\item
  \textbf{Sequence.} \(\hat{P}_1 \circ \hat{P}_2\) executes
  \(\hat{P}_1\) to completion, then \(\hat{P}_2\); both simulations
  compose by the induction hypothesis.
\end{itemize}

The case arguments above are the substantive reduction steps; expanding
each into a fully formal induction over the small-step rules (matching
each WHILE step against the transition functions \(\delta \in \Delta\)
of the corresponding AgentISA instructions) is mechanical and follows
the template of textbook simulation arguments {[}33, 43{]}. We mark
those expansions as routine and leave them to the reader.
\(\blacksquare\)

\subsection{Conclusion}\label{conclusion}

By Lemma A.1, every WHILE-computable function is AgentISA-computable.
Since WHILE is Turing equivalent---a classical result of computability
theory, with WHILE serving as the standard imperative model in semantics
texts such as {[}43{]}---AgentISA is Turing complete. \(\blacksquare\)

\subsection{Scope Remark}\label{scope-remark}

Turing completeness is a deliberately weak property---any language with
unbounded storage, branching, and unbounded iteration attains it. Its
role here is negative certification: the instruction set does not
artificially constrain what agent programs can express, so no
expressiveness objection can be raised against the architecture. The
claim says nothing about efficiency, and it is established without the
LLM-interaction class \(\mathcal{L}\), so it is conservative with
respect to the full instruction set of Definition 3.1.

\section{Proof of Theorem 3.2 (Type
Safety)}\label{proof-of-theorem-3.2-type-safety}

We prove type safety for the pure core by progress and preservation, in
the style of Pierce {[}36{]}. The proof runs over a \textbf{merged
judgment system}: the term typing judgment \(\vdash e : T\) for
\(\lambda\)-terms, and the instruction judgment
\(\vdash i\ \mathit{ok}[pre(i), post(i)]\) for instructions and their
compositions (both defined in B.1). Progress and preservation are each
proved by induction over the derivations of this merged system: the
\(\lambda\)-cases are the standard TAPL argument on term derivations;
the instruction cases are discharged by the standing totality assumption
and the per-operator premises of B.1(ii), with a joint runtime invariant
threaded through the composition cases.

\subsection{Core Calculus}\label{core-calculus}

Restrict attention to the pure core of the instruction language:
expressions over base types, dependent \(\Pi/\Sigma\) pairs, recursion
(\(\mu\)), \(\lambda\)-abstraction and application (whose
\(\beta\)-reduction rule appears in Section 3.1.2), and the composition
operators of Definition 3.6 applied to instructions carrying typed
pre/post-conditions. The small-step relation \(\rightarrow\) is
generated by \(\beta\)-reduction, the context rule
\(E[e_1] \rightarrow E[e_1']\), and, for each instruction \(i\), its
\textbf{fixed per-instruction} transition function \(\delta_i\).
Effectful instructions are treated in B.7.

The two judgment forms:

\begin{enumerate}
\def\labelenumi{(\roman{enumi})}
\item
  \textbf{Term typing} \(\vdash e : T\)---the usual rules over base
  types, extended with the dependent \(\Pi/\Sigma\) introduction forms
  of Definition 3.2.
\item
  \textbf{Instruction typing}
  \(\vdash i\ \mathit{ok}[pre(i), post(i)]\)---read ``\(i\) is a
  well-contracted state transformer'': \(pre(i)\) and \(post(i)\) are
  well-typed predicates over the state type, and the Definition 3.5
  obligation \(\forall s \models pre(i): \delta_i(s) \models post(i)\)
  is discharged per instruction, offline, by the bounded verification of
  Section 3.3.3. Instructions are \textbf{not values and never stuck}:
  totality on precondition-satisfying states is a standing
  assumption---for each instruction \(i\) with its transition function
  \(\delta_i\), \(\delta_i(s)\) is defined whenever
  \(s \models pre(i)\).
\end{enumerate}

Composition rules, for instructions \(i_1, i_2\) carrying contracts:

\[\frac{\vdash i_1\ \mathit{ok} \qquad \vdash i_2\ \mathit{ok} \qquad post(i_1) \Rightarrow pre(i_2)}{\vdash i_1 \circ i_2\ \mathit{ok}[pre(i_1),\, post(i_2)]}\]

\[\frac{\vdash i_1\ \mathit{ok} \qquad \vdash i_2\ \mathit{ok} \qquad pre(i_1) \Leftrightarrow pre(i_2) \quad \text{disjoint state footprints}}{\vdash i_1 \parallel i_2\ \mathit{ok}[pre(i_1),\, post(i_1) \wedge post(i_2)]}\]

\[\frac{\vdash i_1\ \mathit{ok} \qquad \vdash i_2\ \mathit{ok} \qquad post(i_1)\big|_{\text{trigger}} \Rightarrow pre(i_2)}{\vdash i_1 \triangleright i_2\ \mathit{ok}[pre(i_1),\, post(i_1) \wedge (\text{trigger} \Rightarrow post(i_2))]}\]

\[\frac{\vdash i\ \mathit{ok} \qquad \text{guard closed under } post(i) \qquad \text{guard} \Rightarrow pre(i)}{\vdash i^*\ \mathit{ok}[\text{guard},\, post(i)]}\]

Two premises deserve note. The \(\triangleright\) premise is evaluated
\textbf{on the outcome of \(i_1\) that triggers \(i_2\)}: \(i_1\)'s
execution can otherwise destroy \(pre(i_2)\), so the implication the
runtime invariant needs is \(post(i_1)\) restricted to the triggering
outcome, not \(pre(i_1) \Rightarrow pre(i_2)\) on the initial state. The
\(i^*\) premise adds \(\text{guard} \Rightarrow pre(i)\) because the
loop body re-enters: on every re-iteration the guard holds, so the
body's precondition must hold whenever the loop has not terminated.
These rules are generated from---do not add to---the semantics of
Definition 3.6.

\subsection{Canonical Forms}\label{canonical-forms}

If \(\vdash v : T\) and \(v\) is a value, then:

\begin{itemize}
\tightlist
\item
  \(T \in \{Bool, Int, String, Float\}\) implies \(v\) is a literal of
  that type;
\item
  \(T = \Pi(x:T_1).T_2\) implies \(v = \lambda x{:}T_1.\,e\) with
  \(\Gamma, x{:}T_1 \vdash e : T_2\);
\item
  \(T = \Sigma(x:T_1).T_2\) implies \(v = (v_1, v_2)\) with the
  components typed at the corresponding dependent types;
\item
  \(T = \mu X.T\) implies \(v\) is a rolled recursive value;
\item
  no instruction or instruction composition is a value (B.1(ii)):
  instructions are contracted state transformers, not term-level values.
\end{itemize}

\emph{Proof.} By inspection of the typing rules: each type has exactly
one value form, and the syntax of values is defined per type, so the
case analysis is exhaustive. This is the standard canonical-forms lemma
{[}36{]}. \(\blacksquare\)

\subsection{Lemma (Substitution)}\label{lemma-substitution}

If \(\Gamma, x{:}T' \vdash e : T\) and \(\vdash v : T'\), then
\(\Gamma \vdash e[x := v] : T\).

\emph{Proof.} By induction on the derivation of
\(\Gamma, x{:}T' \vdash e : T\). The non-dependent cases are the case
analysis of Pierce {[}36{]}, which we do not duplicate. The dependent
\(\Pi/\Sigma\) cases additionally require synchronizing the substitution
into the second component's type (\(T_2[x := v]\)), as in the
\(\beta\)-case of B.5---the standard dependent-type substitution
bookkeeping, not covered by the non-dependent treatment in {[}36{]}.
\(\blacksquare\)

\subsection{Lemma (Progress)}\label{lemma-progress}

If \(\vdash e : T\), then \(e\) is a value or there exists \(e'\) with
\(e \rightarrow e'\).

\emph{Proof.} By induction on the derivation. The \(\lambda\)-term
cases:

\begin{itemize}
\tightlist
\item
  \textbf{Literals and \(\lambda\)} are values.
\item
  \textbf{Application} \(e_1\,e_2\) with \(e_1 : \Pi(x:T_1).T_2\) and
  \(e_2 : T_1\): by the induction hypothesis, either \(e_1\) steps (then
  the context rule applies) or \(e_1\) is a value, in which case B.2
  gives \(e_1 = \lambda x.\,e\); similarly \(e_2\) either steps or is a
  value; when both are values, \(\beta\)-reduction applies. There is no
  stuck case, because the \(\Pi\)-type of \(e_1\) and B.2 jointly
  guarantee a \(\lambda\) in value position.
\end{itemize}

The instruction cases admit no stuck state by construction. Instructions
step by the standing totality assumption of B.1(ii) whenever their
precondition holds. Compositions maintain the runtime invariant
\textbf{each reached component's precondition holds}, per operator:

\begin{itemize}
\tightlist
\item
  \textbf{Sequencing \(\circ\):} \(pre(i_1)\) holds initially (the
  composition's contract); after \(i_1\) executes,
  \(post(i_1) \Rightarrow pre(i_2)\) re-establishes the invariant for
  \(i_2\).
\item
  \textbf{Parallel \(\parallel\):} disjoint state footprints keep each
  component's precondition valid throughout the interleaving---neither
  component writes what the other's precondition or postcondition reads.
\item
  \textbf{Guarded \(\triangleright\):} on the triggering outcome of
  \(i_1\), the strengthened premise
  \(post(i_1)\big|_{\text{trigger}} \Rightarrow pre(i_2)\)
  re-establishes \(pre(i_2)\); on the non-triggering outcome, \(i_2\)
  does not run.
\item
  \textbf{Iteration \(i^*\):} the body re-enters only while the guard
  holds, and the premise \(\text{guard} \Rightarrow pre(i)\) makes each
  re-entry legitimate; when the guard fails, the loop is terminal.
\end{itemize}

The invariant is threaded jointly with B.5: each individual step
preserves the invariant by the per-operator argument above. For
iteration, either the guard is false (terminal) or the body's
precondition holds by \(\text{guard} \Rightarrow pre(i)\), and the body
steps by the induction hypothesis on the component derivation.
\(\blacksquare\)

\subsection{Lemma (Preservation)}\label{lemma-preservation}

If \(\vdash e : T\) and \(e \rightarrow e'\), then \(\vdash e' : T\).

\emph{Proof.} By induction on the derivation, with case analysis on the
step:

\begin{itemize}
\tightlist
\item
  \textbf{\(\beta\):} \((\lambda x.\,e)\,v \rightarrow e[x := v]\). The
  application rule types the redex at \(T_2[x := v]\); by B.3, the
  substituted term \(e[x := v]\) types at \(T_2[x := v]\), which is
  exactly the type the application rule assigned (the dependent-type
  bookkeeping; when \(x\) does not occur in \(T_2\) this degenerates to
  the usual \(T_2 = T\)).
\item
  \textbf{Context:} \(E[e_1] \rightarrow E[e_1']\) with
  \(e_1 \rightarrow e_1'\). By the induction hypothesis \(e_1' : T_1\);
  an induction on the structure of \(E\) re-types \(E[e_1']\) at the
  original type.
\item
  \textbf{Instruction transitions:} the instruction's \(\mathit{ok}\)
  judgment (B.1(ii)) certifies that \(\delta_i\) establishes \(post(i)\)
  from any \(pre(i)\)-satisfying state (discharged offline by the
  bounded verification of Section 3.3.3), so the reduct re-types at the
  postcondition.
\item
  \textbf{Composition operators:} each operator's step is typed by its
  B.1(ii) rule; the reduct's contract follows from the operator's stated
  conclusion (for \(\circ\), \(post(i_2)\); for \(\parallel\), the
  conjunction; for \(\triangleright\), the trigger-restricted
  conjunction; for \(i^*\), \(post(i)\)), by the induction hypothesis on
  the component derivations.
\end{itemize}

The remaining instruction-class cases repeat this pattern with
class-specific invariants; they are routine and omitted (the template is
the standard preservation argument for the corresponding constructs
{[}36{]}). \(\blacksquare\)

\subsection{Theorem 3.2}\label{theorem-3.2}

By B.4, a well-typed closed term is a value or takes a step; by B.5,
every step preserves typing. Hence no well-typed term reaches a stuck
state: well-typed programs in the pure core cannot go wrong.
\(\blacksquare\)

\subsection{Scope Remarks}\label{scope-remarks}

Two caveats state the theorem's boundaries honestly. First, B.4's appeal
to totality of transition functions on precondition-satisfying states is
an assumption about instruction \emph{implementations}, not something
the type system alone discharges. The per-instruction mechanism intended
to discharge it is the bounded verification of Section 3.3.3, which is
now shipped code: the checker mechanically discharges the obligation for
pure instructions on the state spaces it can enumerate---exhaustively
for small enumerable domains, by deterministic bounded sampling
beyond---while side-effectful instructions and state spaces beyond the
enumeration bound remain covered by hand argument and unit tests rather
than by the mechanical check. The theorem is conditional on the
obligation being (or becoming) true of each instruction, to the extent
stated. Second, effectful instructions (LLM calls, resource allocation)
can fail at runtime for reasons outside the type system (network errors,
exhausted capacity); Theorem 3.2 covers the pure core and effectful
instructions with checkable preconditions. Runtime failure handling for
the remaining cases is the responsibility of the fallback machinery of
Section 5.6, not of types.

\section{Proofs of Theorems 3.4--3.6}\label{proofs-of-theorems-3.43.6}

\subsection{Denotational Setup}\label{denotational-setup}

Fix the state space \(\Sigma\) (Definition 3.1's state space---not the
dependent-type constructor \(\Sigma\) of Definition 3.2). For a pure
instruction \(i\), let
\(\llbracket i \rrbracket : \Sigma \rightharpoonup \Sigma\) denote its
denotation, a partial function total on states satisfying \(pre(i)\).
Program denotation is defined by structural recursion on the composition
algebra of Definition 3.6:

\begin{itemize}
\tightlist
\item
  \(\llbracket i \rrbracket\) for atomic instructions (given);
\item
  \(\llbracket p_1 \circ p_2 \rrbracket = \llbracket p_2 \rrbracket \circ \llbracket p_1 \rrbracket\)
  (note the order: \(p_1\) executes first);
\item
  \(\llbracket p_1 \parallel p_2 \rrbracket(\sigma) = \llbracket p_2 \rrbracket(\llbracket p_1 \rrbracket(\sigma))\)
  for pure \(p_1, p_2\) with disjoint state footprints (as the
  \(\parallel\)-typing rule requires)---the order is immaterial because
  updates to disjoint cells commute on their domains;
\item
  \(\llbracket p_1 \triangleright p_2 \rrbracket(\sigma) = \llbracket p_2 \rrbracket(\llbracket p_1 \rrbracket(\sigma))\)
  when \(p_1\)'s outcome at \(\sigma\) is the trigger, and
  \(\llbracket p_1 \rrbracket(\sigma)\) otherwise (the denotation is
  defined on \(\Sigma + \text{outcome}\), with \(p_2\) applied
  conditionally on the trigger tag);
\item
  \(\llbracket p^* \rrbracket\) is the least fixed point of the
  corresponding functional.
\end{itemize}

Write \(p \equiv q\) iff
\(\llbracket p \rrbracket = \llbracket q \rrbracket\) on all states
where both are defined.

\subsection{Lemma (Atomic Rewrite
Soundness)}\label{lemma-atomic-rewrite-soundness}

Each rewrite rule implemented by the optimizer (Section 3.4.4) preserves
denotation.

\textbf{(i) Fusion.} \(i_1 \circ i_2 \rightarrow i_{12}\), with
\(pre(i_{12}) = pre(i_1)\), \(post(i_{12}) = post(i_2)\), and
\(i_{12}\)'s body stipulated to compute
\(\llbracket i_2 \rrbracket \circ \llbracket i_1 \rrbracket\)
(intermediate bindings named, so the composition is explicit). The
denotational claim is then an instance of the definition of \(\circ\)
(C.1): a body computing
\(\llbracket i_2 \rrbracket \circ \llbracket i_1 \rrbracket\) denotes
exactly \(\llbracket i_1 \circ i_2 \rrbracket\) (the associativity of
function composition is what fuses \emph{three or more} consecutive
instructions). \textbf{Two honest boundaries.} First, this is an
equality \emph{of the rewrite algebra}: it holds of the fused form whose
body computes the composite. The shipped \emph{formal compiler} (the
static path of the Section 5.3 benchmarks) produces a formal fused
\emph{node} (merged metadata and dependencies) without generating an
executable composite body---code generation for fused nodes is not
implemented---so what those benchmarks measure is the structural fusion
of the formal stream, and the executable-fusion claim would additionally
require the codegen step. The \emph{Instruction-level} merge path (the
dynamic optimizer of Sections 4.2 and 5.4) does carry an executable
body---the \texttt{CompositeInstruction} replay construction---and its
differential property (composite execution equals constituent-sequence
execution; the fusion gate refuses producer--consumer pairs) is pinned
by unit tests and exercised by the trace-replay identity of Section 5.4
(5/5 seeds). A fused \emph{single} executable body distinct from replay
remains specified. Second, the compiler applies fusion only when output
types of \(i_1\) match input types of \(i_2\) and both members are pure,
since denotational equality of the fused form presupposes
order-insensitivity of intermediate effects, and the shipped gate
achieves it via purity (the pure or order-insensitive alternative is
specified, not implemented).

\textbf{(ii) Dead-code elimination.} Instruction \(i\) is removed from
\(p\) when \(Def(i) \cap Live_{out}(i) = \emptyset\), with liveness the
standard fixpoint of
\(Live_{in}(i) = Use(i) \cup (Live_{out}(i) - Def(i))\) {[}33{]}. Claim:
if \(i\) is dead and side-effect-free, then
\(p \equiv p \setminus \{i\}\). Proof sketch: by induction on the suffix
program following \(i\), its denotation depends only on the live
variables at \(i\)'s position; \(i\) writes none of them, so removing
\(i\) cannot change any value subsequently read. The
side-effect-freeness gate matches the purity restriction of Definition
3.6 (Property 2 and the Remark following it)---instructions declared
\texttt{side\_effects=True} are never eliminated.

\textbf{(iii) Peephole combining.} Each rule
\((s_1, \dots, s_k) \rightarrow s'\) carries the verification obligation
\(\forall \sigma: \llbracket s_1 \circ \dots \circ s_k \rrbracket(\sigma) = \llbracket s' \rrbracket(\sigma)\),
to be discharged at rule-definition time (Section 4.1.3). Candidate
instances, \emph{illustrative only---the shipped optimizer registers
none of them}: \texttt{LOAD\ A;\ STORE\ B} composes to
\texttt{MOVE\ A→B} (degenerate to the identity when A = B, since the
value read is the value written; the general rule of Section 3.4.4);
consecutive \texttt{QUERY} instructions with identical (instruction,
params) pairs collapse to one query whose result is reused (the
batch-QUERY rule of Section 3.4.4, specialized to identical pairs) (an
obligation about a \emph{compile-time} rewrite; the caching measured in
Section 5.2 is a different, \emph{runtime} mechanism---a result cache
consulted at execution time---and its numbers do not evidence this
rule); consecutive \texttt{CONTEXT\_UPDATE}s compose to a single update
by associativity of the update function.

\emph{Proof.} Each case reduces to the construction of the rewritten
instruction and the definition of liveness. The inner inductions---on
the pair of bodies for (i), on suffix programs for (ii)---are routine
and left to the reader. \(\blacksquare\)

\subsection{Lemma (Compositionality)}\label{lemma-compositionality}

Let \(T\) be a program transform with
\(\llbracket T(p) \rrbracket = \llbracket p \rrbracket\) for all \(p\).
Then for any context \(C[\cdot]\) built from the composition operators,
\(\llbracket C[T(p)] \rrbracket = \llbracket C[p] \rrbracket\).

\emph{Proof.} By structural induction on \(C\). Each operator's
denotation (C.1) is defined compositionally in its constituents'
denotations, so substituting a denotationally equal constituent leaves
the whole unchanged. The parallel case uses the disjoint-footprint
hypothesis of the \(\parallel\)-typing rule (updates to disjoint cells
commute). The iteration case uses the defining property of least fixed
points: if the functionals induced by \(p\) and \(T(p)\) are equal, then
all finite approximants are equal (by the induction hypothesis), so
their suprema---the least fixed points---are equal. \(\blacksquare\)

\subsection{Theorem 3.4 (Optimization
Soundness)}\label{theorem-3.4-optimization-soundness}

\emph{Proof.} Read all equalities below as the partial-equivalence
\(\equiv\) of C.1 (equality on all states where both sides are defined),
since the DCE rule preserves denotation only on the defined domain and
can strictly enlarge it. The compiler applies a finite sequence of
rewrites \(O = R_k \circ \dots \circ R_1\) (finiteness follows from
compilation terminating; each pass is a bounded traversal). By induction
on \(k\): \(k = 0\) is trivial. For the step, write \(O = R \circ O'\).
By the induction hypothesis
\(\llbracket O'(p) \rrbracket \equiv \llbracket p \rrbracket\). Two
routes, by rule kind. \textbf{Fusion and peephole} are local subprogram
rewrites: the atomic-rewrite lemma (C.2) preserves the denotation of the
subprogram it rewrites, and the compositionality lemma (C.3) lifts the
preservation through any surrounding composition context.
\textbf{Dead-code elimination} is not local---liveness is computed over
the whole surrounding program---so it is handled at pass level directly
by C.2(ii)'s whole-program argument (\(p \equiv p \setminus \{i\}\)),
not through the lifting template. Hence
\(\llbracket R(O'(p)) \rrbracket \equiv \llbracket O'(p) \rrbracket \equiv \llbracket p \rrbracket\).
\(\blacksquare\)

\subsection{Theorem 3.5 (Pipeline
Correctness)}\label{theorem-3.5-pipeline-correctness}

Let \(s\) be a valid instruction sequence and \(S\) a schedule
satisfying:

\begin{itemize}
\tightlist
\item
  \textbf{(S1)} every instruction is dispatched exactly once;
\item
  \textbf{(S2)} for every dependency \((i_1, i_2) \in Deps\):
  \(start(i_2) \geq end(i_1)\) (the constraint of Section 3.4.3);
\item
  \textbf{(S3)} instructions sharing a parallel dispatch group are
  pairwise non-interfering (no data hazard, no resource conflict);
\item
  \textbf{(S4)} commit is in program order.
\end{itemize}

Then executing \(s\) under \(S\) yields the same final state as
sequential execution:
\(\llbracket s \rrbracket_S = \llbracket s \rrbracket_{seq}\).

\emph{Proof.} Let \(g_1, \dots, g_m\) be the dispatch groups of \(S\) in
commit order. By S4, groups commit in program order and group commits do
not interleave, so the union \(\bigcup_{j \leq t} g_j\) of committed
groups is always a program-order prefix---``executing all instructions
in those groups'' below means executing that prefix sequentially. We
show the invariant

\[\begin{gathered} I(t): \text{after committing } g_1, \dots, g_t,\\ \text{the architectural state equals the sequential state after}\\ \text{executing all instructions in those groups.}\end{gathered}\]

\begin{itemize}
\tightlist
\item
  \textbf{Base} (\(t = 0\)): both sides are the initial state.
\item
  \textbf{Step}: assume \(I(t)\). Every instruction in \(g_{t+1}\) has
  its dependencies in earlier groups (by S2, since it is scheduled after
  them), which by \(I(t)\) have completed and committed. Within the
  group, S3 gives pairwise non-interference: no instruction reads
  another group-member's output, and no two write the same location or
  resource. Hence the group's updates commute---for pure updates this is
  the diamond property, and for effectful instructions, S3's exclusion
  of shared resources makes their effects independent, so the combined
  effect is unique regardless of interleaving. The state after
  \(g_{t+1}\) is therefore independent of the group's internal
  interleaving and equals the sequential execution of the same
  instructions from the state at \(I(t)\).
\item
  By induction, \(I(m)\) holds; by S1 no instruction was skipped or
  duplicated.
\end{itemize}

This completes the induction over all dispatch groups, hence the proof.
\(\blacksquare\)

\textbf{Corollary (dispatch folding).} The ADVANCED and AGGRESSIVE
levels fold instructions into shared dispatches (Section 5.3). A folded
group satisfies S3 by construction: the grouping check admits only
side-effect-free instructions whose dependencies and declared outputs
are pairwise disjoint, so the group's members act on disjoint state
components and their updates commute (Definition 3.6, Property 2). By
Theorem 3.5 the fold changes the \emph{dispatch count} but neither the
semantics nor---consistent with the measurement of Section 5.3---the
effective instruction count.

\subsection{Theorem 3.6 (Pass-Pipeline
Soundness)}\label{theorem-3.6-pass-pipeline-soundness}

Every optimization in the pass pipeline is a finite composition of
rewrites covered by the atomic-rewrite lemma (C.2) or by the
dispatch-folding corollary of Theorem 3.5 (C.5)---the grouping pass is a
fold, not a local rewrite---and the argument of C.4 applies verbatim to
each rewrite layer. \(\blacksquare\)

\subsection{Scope Remark}\label{scope-remark-1}

Theorems 3.4--3.6 are \emph{conditional}: they hold of transformations
whose verification obligations (Sections 4.1.3 and 4.3.1) are
discharged. As stated there, those obligations are \textbf{specified but
not yet mechanically discharged}---the shipped optimizer registers no
peephole rules (the peephole obligation is dormant, and the
fusion/DCE/grouping passes are argued directly in C.2 and C.5 rather
than through the symbolic-execution checker), and the symbolic-execution
equivalence checker itself is a specification (Section 4.3.1). Where an
obligation is discharged only on bounded state spaces (Section 3.3.3),
the guarantee is correspondingly bounded: rewrites verified on the
explored subset are sound there. No such rewrite is registered today;
when one is, its application must be restricted to the explored subset
(that restriction is a dormant obligation of the future rule set, not
shipped behavior). What the shipped implementation does provide is
empirical semantics-preservation at task level: the system-level
demonstration of Section 5.8 re-executes an optimized stream and checks
byte-identical output.

\textbf{Observation model.} The theorems' equality is \emph{partial
equivalence} (C.1: \(\llbracket p \rrbracket = \llbracket q \rrbracket\)
on states where both are defined), and the observation model should be
read accordingly. \textbf{Observed:} the final state of declared state
cells, the declared outputs (the instruction's product keys), and---for
theorems covering effectful instructions (C.5)---the externally visible
effects, in program order. \textbf{Deliberately not observed:} the
\emph{domain of definition} itself. Because equivalence is only partial,
a rewrite may extend definedness---DCE is the canonical example (C.4):
deleting an instruction that would have failed can turn a previously
undefined run into a defined one. If program failure or exception
behavior were part of the observation model, this would violate
preservation; it is therefore stated explicitly that failure behavior is
outside what Theorems 3.4--3.6 preserve, and any deployment for which
failures are observable (e.g., a pipeline contract that treats
exceptions as data) needs a strengthened equivalence and must re-derive
the theorems under it. Internal traces---the sequence of instructions
executed---are likewise not observed; only the net transformation is.

\textbf{Three status levels.} Each mechanism in this paper falls into
exactly one of three levels, which we keep separate throughout:
\emph{(i) specified}---the symbolic-execution equivalence checker, the
pipeline scheduler's cycle-accurate extension, tree-pattern instruction
selection (Section 4.1.2); \emph{(ii) implemented with checks}---the
optimizer passes (fusion, DCE, grouping) with their pass-level arguments
in C.2/C.5, the runtime monitors, and differential unit tests pinning
the fusion property this appendix argues (fused-composite execution
equals constituent-sequence execution; the fusion gate refuses
producer--consumer pairs;
\texttt{tests/core/isa/test\_fusion\_differential.py}); \emph{(iii)
empirically validated}---trace-level replay identity on the mock
(Section 5.4: optimized traces re-executed, final states identical in
5/5 seeds), task-level re-execution with byte-identical outputs (Section
5.8), and the caching/compilation measurements of Section 5 (including
the exact-response-cache parity baseline of Section 5.2). A statement at
one level is never used to imply a guarantee at a stronger level.

\section{Specified-Only Component
Listings}\label{specified-only-component-listings}

This appendix carries the algorithm listings that the main paper
specifies but does not ship (or ships dormant), moved from its Section 4
to respect the page limit. Each listing keeps its original framing;
implementation status statements are unchanged. The shipped selection
step (the heuristic name-overlap matcher) remains in the source module
\texttt{agentflow.core.isa.instruction\_selector}.

\subsection{Specified Tree-Pattern Instruction Selection
(MATCH-TREE)}\label{specified-tree-pattern-instruction-selection-match-tree}

Input: IR tree T, pattern set P; output: optimal pattern selection S.

\begin{Shaded}
\begin{Highlighting}[]
\NormalTok{Input: IR tree T, Pattern }\BuiltInTok{set}\NormalTok{ P}
\NormalTok{Output: Optimal pattern selection S}

\NormalTok{function MATCH}\OperatorTok{{-}}\NormalTok{TREE(T, P)}
\NormalTok{    costs[n] ← ∞ }\ControlFlowTok{for} \BuiltInTok{all}\NormalTok{ nodes n }\KeywordTok{in}\NormalTok{ T}
\NormalTok{    Let selected[n] }\OperatorTok{=}\NormalTok{ ∅ }\ControlFlowTok{for} \BuiltInTok{all}\NormalTok{ nodes n }\KeywordTok{in}\NormalTok{ T}
    
\NormalTok{    function MATCH}\OperatorTok{{-}}\NormalTok{SUBTREE(node)}
        \ControlFlowTok{if}\NormalTok{ costs[node] ≠ ∞ then}
            \ControlFlowTok{return}\NormalTok{ costs[node]}
\NormalTok{        end }\ControlFlowTok{if}
        
\NormalTok{        min\_cost ← ∞}
\NormalTok{        best\_pat ← null}
        
        \ControlFlowTok{for}\NormalTok{ each pattern p }\KeywordTok{in}\NormalTok{ P do}
            \ControlFlowTok{if}\NormalTok{ MATCHES(p, node) then}
\NormalTok{                curr\_cost ← PATTERN}\OperatorTok{{-}}\NormalTok{COST(p)}
                \ControlFlowTok{for}\NormalTok{ each child c of node }\KeywordTok{not}\NormalTok{ covered by p do}
\NormalTok{                    curr\_cost ← curr\_cost }\OperatorTok{+}\NormalTok{ MATCH}\OperatorTok{{-}}\NormalTok{SUBTREE(c)}
\NormalTok{                end }\ControlFlowTok{for}
                
                \ControlFlowTok{if}\NormalTok{ curr\_cost }\OperatorTok{\textless{}}\NormalTok{ min\_cost then}
\NormalTok{                    min\_cost ← curr\_cost}
\NormalTok{                    best\_pat ← p}
\NormalTok{                end }\ControlFlowTok{if}
\NormalTok{            end }\ControlFlowTok{if}
\NormalTok{        end }\ControlFlowTok{for}
        
\NormalTok{        costs[node] ← min\_cost}
\NormalTok{        selected[node] ← best\_pat}
        \ControlFlowTok{return}\NormalTok{ min\_cost}
\NormalTok{    end function}
    
\NormalTok{    MATCH}\OperatorTok{{-}}\NormalTok{SUBTREE(root(T))}
    \ControlFlowTok{return}\NormalTok{ (costs, selected)}
\NormalTok{end function}
\end{Highlighting}
\end{Shaded}

\textbf{Theorem 4.1} (Selection Optimality, specified algorithm): given
a set of context-local patterns P and an IR tree T, the MATCH-TREE
algorithm computes the minimal covering cost; the covering itself is
read off the \texttt{selected} table afterwards, by walking from the
root and skipping subtrees spanned by the chosen pattern.

\emph{Proof}: By induction on tree height, assuming a covering of T
exists (if no pattern matches a node its cost is \(\infty\) and no
covering exists):

\begin{enumerate}
\def\labelenumi{\arabic{enumi}.}
\item
  Base case (a leaf has no children): the recurrence degenerates to a
  plain argmin over the patterns matching the leaf, which is optimal for
  a single node.
\item
  Inductive step: For a node \(v\) with children \(c_1, \dots, c_k\): \[
  \text{cost}(v) = \min_{p \in P} \left\{p.\text{cost} + \sum_{\substack{c \in \text{children}(v) \\ c \notin \text{covered}(p)}} \text{cost}(c)\right\}
  \]

  By the induction hypothesis, each \(\text{cost}(c)\) is the minimal
  covering cost of the corresponding subtree; the recurrence takes each
  candidate pattern's cost plus the (already optimal) costs of the
  children the candidate does not span, so the minimum over candidate
  patterns is the minimal covering cost at \(v\)---the minimum over
  disjoint pattern choices is optimal.
\end{enumerate}

The recurrence sums only over children not covered by the candidate
pattern \(p\)---summing over all children would double-count the cost of
subtrees the pattern spans. Section 3.4.3's per-node argmin is therefore
exact for single-node patterns; multi-node patterns require this
covering recurrence.

The recurrence requires that a pattern's feasibility at a node not
depend on the covers chosen for other subtrees; context constraints that
couple sibling subtrees break this optimal-substructure property, which
is why the statement restricts to context-local patterns. This theorem
is a property of the specified algorithm, not of the shipped heuristic
selector; no optimality claim is made for the implementation.

\subsection{Specified Cycle-Accurate Pipeline Scheduler
(SCHEDULE-PIPELINE)}\label{specified-cycle-accurate-pipeline-scheduler-schedule-pipeline}

Input: instruction sequence I; output: schedule S with cycle
assignments.

\begin{Shaded}
\begin{Highlighting}[]
\NormalTok{Input: Instruction sequence I}
\NormalTok{Output: Schedule S }\ControlFlowTok{with}\NormalTok{ cycle assignments}

\NormalTok{function SCHEDULE}\OperatorTok{{-}}\NormalTok{PIPELINE(I)}
\NormalTok{    D ← BUILD}\OperatorTok{{-}}\NormalTok{DEPENDENCY}\OperatorTok{{-}}\NormalTok{GRAPH(I)}
\NormalTok{    S ← empty schedule}
\NormalTok{    R ← priority queue of ready instructions}
    
    \ControlFlowTok{while}\NormalTok{ D }\KeywordTok{not}\NormalTok{ empty do}
\NormalTok{        c ← CURRENT}\OperatorTok{{-}}\NormalTok{CYCLE()}
        
        \OperatorTok{//}\NormalTok{ Find newly ready instructions}
\NormalTok{        N ← GET}\OperatorTok{{-}}\NormalTok{READY}\OperatorTok{{-}}\NormalTok{NODES(D, c)}
\NormalTok{        ADD}\OperatorTok{{-}}\NormalTok{TO}\OperatorTok{{-}}\NormalTok{QUEUE(R, N)}
        
        \OperatorTok{//}\NormalTok{ Schedule instructions }\ControlFlowTok{for}\NormalTok{ current cycle}
        \ControlFlowTok{while}\NormalTok{ R }\KeywordTok{not}\NormalTok{ empty }\KeywordTok{and}\NormalTok{ HAS}\OperatorTok{{-}}\NormalTok{RESOURCES(c) do}
\NormalTok{            i ← POP}\OperatorTok{{-}}\NormalTok{HIGHEST}\OperatorTok{{-}}\NormalTok{PRIORITY(R)}
\NormalTok{            slot ← ALLOCATE}\OperatorTok{{-}}\NormalTok{RESOURCES(i, c)}
\NormalTok{            ADD}\OperatorTok{{-}}\NormalTok{TO}\OperatorTok{{-}}\NormalTok{SCHEDULE(S, i, slot)}
\NormalTok{            UPDATE}\OperatorTok{{-}}\NormalTok{DEPENDENCIES(D, i)}
\NormalTok{        end }\ControlFlowTok{while}
\NormalTok{    end }\ControlFlowTok{while}
    
\NormalTok{    VERIFY}\OperatorTok{{-}}\NormalTok{SCHEDULE(S)}
    \ControlFlowTok{return}\NormalTok{ S}
\NormalTok{end function}
\end{Highlighting}
\end{Shaded}

\subsection{Specified Dynamic Pipeline Controller
(PipelineController)}\label{specified-dynamic-pipeline-controller-pipelinecontroller}

\begin{Shaded}
\begin{Highlighting}[]
\KeywordTok{class}\NormalTok{ PipelineController:}
    \KeywordTok{def} \FunctionTok{\_\_init\_\_}\NormalTok{(}\VariableTok{self}\NormalTok{):}
        \VariableTok{self}\NormalTok{.hazard\_detector }\OperatorTok{=}\NormalTok{ HazardDetector()}
        \VariableTok{self}\NormalTok{.resource\_monitor }\OperatorTok{=}\NormalTok{ ResourceMonitor()}
        \VariableTok{self}\NormalTok{.performance\_tracker }\OperatorTok{=}\NormalTok{ PerformanceTracker()}
        
    \ControlFlowTok{async} \KeywordTok{def}\NormalTok{ execute\_pipeline(}\VariableTok{self}\NormalTok{, }
\NormalTok{                             instruction\_stream: AsyncIterator[Instruction]}
\NormalTok{                            ) }\OperatorTok{{-}\textgreater{}}\NormalTok{ AsyncIterator[Result]:}
        \ControlFlowTok{async} \ControlFlowTok{for}\NormalTok{ inst }\KeywordTok{in}\NormalTok{ instruction\_stream:}
            \CommentTok{\# Hazard detection}
            \ControlFlowTok{while} \VariableTok{self}\NormalTok{.hazard\_detector.has\_hazard(inst):}
                \ControlFlowTok{await} \VariableTok{self}\NormalTok{.hazard\_detector.wait\_resolution()}
                
            \CommentTok{\# Resource allocation}
\NormalTok{            resources }\OperatorTok{=} \ControlFlowTok{await} \VariableTok{self}\NormalTok{.resource\_monitor.allocate(inst)}
            
            \ControlFlowTok{try}\NormalTok{:}
                \CommentTok{\# Execution}
\NormalTok{                result }\OperatorTok{=} \ControlFlowTok{await} \VariableTok{self}\NormalTok{.execute\_instruction(inst, resources)}
                \CommentTok{\# Performance tracking}
                \VariableTok{self}\NormalTok{.performance\_tracker.record(inst, result)}
                \ControlFlowTok{yield}\NormalTok{ result}
            \ControlFlowTok{finally}\NormalTok{:}
                \VariableTok{self}\NormalTok{.resource\_monitor.release(resources)}
\end{Highlighting}
\end{Shaded}

\subsection{Specified Hot-Path Detector
(DETECT-HOT-PATH)}\label{specified-hot-path-detector-detect-hot-path}

Input: execution path P, threshold \(\theta\), EMA decay \(\gamma\)
(\(\gamma\) = 2/(w+1) for window w; \(\gamma\) is unrelated to the
\(\alpha\) stability bound of the main paper's Section 3.4.4); output:
Boolean indicating if P is hot.

\begin{Shaded}
\begin{Highlighting}[]
\NormalTok{Input: Execution path P, Threshold θ, EMA decay γ (γ }\OperatorTok{=} \DecValTok{2}\OperatorTok{/}\NormalTok{(w}\OperatorTok{+}\DecValTok{1}\NormalTok{) }\ControlFlowTok{for}\NormalTok{ window w}\OperatorTok{;}\NormalTok{ γ }\KeywordTok{is}\NormalTok{ unrelated to the α stability bound of Section }\FloatTok{3.4.4}\NormalTok{)}
\NormalTok{Output: Boolean indicating }\ControlFlowTok{if}\NormalTok{ P }\KeywordTok{is}\NormalTok{ hot}

\NormalTok{function DETECT}\OperatorTok{{-}}\NormalTok{HOT}\OperatorTok{{-}}\NormalTok{PATH(P, θ, γ)}
\NormalTok{    Let counter be an exponential moving average }\ControlFlowTok{with}\NormalTok{ decay γ}
    
\NormalTok{    function UPDATE}\OperatorTok{{-}}\NormalTok{COUNTERS(path)}
\NormalTok{        h ← COMPUTE}\OperatorTok{{-}}\NormalTok{PATH}\OperatorTok{{-}}\NormalTok{HASH(path)}
\NormalTok{        counter[h] ← counter[h] × (}\DecValTok{1} \OperatorTok{{-}}\NormalTok{ γ) }\OperatorTok{+} \DecValTok{1}
        
        \ControlFlowTok{if}\NormalTok{ IS}\OperatorTok{{-}}\NormalTok{PATH}\OperatorTok{{-}}\NormalTok{HOT(path) then}
\NormalTok{            TRIGGER}\OperatorTok{{-}}\NormalTok{OPTIMIZATION(path)}
\NormalTok{        end }\ControlFlowTok{if}
\NormalTok{    end function}
    
\NormalTok{    function IS}\OperatorTok{{-}}\NormalTok{PATH}\OperatorTok{{-}}\NormalTok{HOT(path)}
\NormalTok{        h ← COMPUTE}\OperatorTok{{-}}\NormalTok{PATH}\OperatorTok{{-}}\NormalTok{HASH(path)}
\NormalTok{        freq ← counter[h] }\OperatorTok{/}\NormalTok{ Σ\_\{h}\StringTok{\textquotesingle{}\} counter[h\textquotesingle{}}\NormalTok{]   }\OperatorTok{//}\NormalTok{ recent execution share, }\KeywordTok{not}\NormalTok{ cumulative total}
\NormalTok{        stab ← COMPUTE}\OperatorTok{{-}}\NormalTok{STABILITY(path)}
        \ControlFlowTok{return}\NormalTok{ freq }\OperatorTok{\textgreater{}}\NormalTok{ θ }\KeywordTok{and}\NormalTok{ stab }\OperatorTok{\textgreater{}}\NormalTok{ STABILITY}\OperatorTok{{-}}\NormalTok{THRESHOLD}
\NormalTok{    end function}
    
\NormalTok{    UPDATE}\OperatorTok{{-}}\NormalTok{COUNTERS(P)}
    \ControlFlowTok{return}\NormalTok{ IS}\OperatorTok{{-}}\NormalTok{PATH}\OperatorTok{{-}}\NormalTok{HOT(P)}
\NormalTok{end function}
\end{Highlighting}
\end{Shaded}

\subsection{Specified Trace Formation
(BUILD-TRACE)}\label{specified-trace-formation-build-trace}

Input: hot path H; output: optimized trace T.

\begin{Shaded}
\begin{Highlighting}[]
\NormalTok{Input: Hot path H}
\NormalTok{Output: Optimized trace T}

\NormalTok{function BUILD}\OperatorTok{{-}}\NormalTok{TRACE(H)}
\NormalTok{    T ← empty trace}
    
\NormalTok{    function SHOULD}\OperatorTok{{-}}\NormalTok{INCLUDE(block)}
        \ControlFlowTok{return}\NormalTok{ EXECUTION}\OperatorTok{{-}}\NormalTok{FREQUENCY(block) }\OperatorTok{\textgreater{}}\NormalTok{ THRESHOLD }\KeywordTok{and}
\NormalTok{               PREDICTABLE}\OperatorTok{{-}}\NormalTok{BRANCHES(block)}
\NormalTok{    end function}
    
\NormalTok{    function SHOULD}\OperatorTok{{-}}\NormalTok{SPLIT(block)}
        \ControlFlowTok{return}\NormalTok{ HAS}\OperatorTok{{-}}\NormalTok{SIDE}\OperatorTok{{-}}\NormalTok{EXIT(block) }\KeywordTok{or}
\NormalTok{               EXECUTION}\OperatorTok{{-}}\NormalTok{FREQUENCY(block) }\OperatorTok{\textless{}}\NormalTok{ MIN}\OperatorTok{{-}}\NormalTok{FREQ}
\NormalTok{    end function}
    
    \ControlFlowTok{for}\NormalTok{ each block B }\KeywordTok{in}\NormalTok{ H do}
        \ControlFlowTok{if}\NormalTok{ SHOULD}\OperatorTok{{-}}\NormalTok{INCLUDE(B) then}
\NormalTok{            APPEND}\OperatorTok{{-}}\NormalTok{TO}\OperatorTok{{-}}\NormalTok{TRACE(T, B)}
        \ControlFlowTok{else}
            \ControlFlowTok{break}   \OperatorTok{//} \BuiltInTok{any}\NormalTok{ block }\KeywordTok{not}\NormalTok{ selected terminates the trace}
\NormalTok{        end }\ControlFlowTok{if}
\NormalTok{    end }\ControlFlowTok{for}
    
\NormalTok{    VERIFY}\OperatorTok{{-}}\NormalTok{TRACE}\OperatorTok{{-}}\NormalTok{PROPERTIES(T)}
    \ControlFlowTok{return}\NormalTok{ T}
\NormalTok{end function}
\end{Highlighting}
\end{Shaded}

\subsection{Specified Symbolic-Equivalence Checker
(VERIFY-TRANSFORMATION)}\label{specified-symbolic-equivalence-checker-verify-transformation}

Input: original CFG O, transformed CFG T; output: Boolean indicating
semantic equivalence.

\begin{Shaded}
\begin{Highlighting}[]
\NormalTok{Input: Original CFG O, Transformed CFG T}
\NormalTok{Output: Boolean indicating semantic equivalence}

\NormalTok{function VERIFY}\OperatorTok{{-}}\NormalTok{TRANSFORMATION(O, T)}
    \OperatorTok{//}\NormalTok{ Build symbolic execution paths}
\NormalTok{    P₁ ← SYMBOLIC}\OperatorTok{{-}}\NormalTok{EXECUTE(O)}
\NormalTok{    P₂ ← SYMBOLIC}\OperatorTok{{-}}\NormalTok{EXECUTE(T)}
    
\NormalTok{    function VERIFY}\OperatorTok{{-}}\NormalTok{EQUIVALENCE(path₁, path₂)}
        \OperatorTok{//}\NormalTok{ σ ranges over the symbolic states Σ explored on both paths}
        \ControlFlowTok{for}\NormalTok{ each symbolic state σ ∈ Σ do}
            \ControlFlowTok{if}\NormalTok{ EVALUATE(path₁, σ) ≠ EVALUATE(path₂, σ) then}
                \ControlFlowTok{return}\NormalTok{ false}
\NormalTok{            end }\ControlFlowTok{if}
\NormalTok{        end }\ControlFlowTok{for}
        \ControlFlowTok{return}\NormalTok{ true}
\NormalTok{    end function}
    
    \OperatorTok{//}\NormalTok{ Paths must pair one}\OperatorTok{{-}}\NormalTok{to}\OperatorTok{{-}}\NormalTok{one}\OperatorTok{;}\NormalTok{ a length mismatch means the}
    \OperatorTok{//}\NormalTok{ transformation added }\KeywordTok{or}\NormalTok{ removed a symbolic path — reject.}
    \ControlFlowTok{if}\NormalTok{ LENGTH(P₁) ≠ LENGTH(P₂) then}
        \ControlFlowTok{return}\NormalTok{ false}
\NormalTok{    end }\ControlFlowTok{if}
    
    \OperatorTok{//}\NormalTok{ Verify }\BuiltInTok{all}\NormalTok{ path pairs, paired by the transformation}\StringTok{\textquotesingle{}s path map}
\ErrorTok{    for each }\NormalTok{(p₁, p₂) }\KeywordTok{in}\NormalTok{ PATH}\OperatorTok{{-}}\NormalTok{MAP(P₁, P₂) do}
        \ControlFlowTok{if} \KeywordTok{not}\NormalTok{ VERIFY}\OperatorTok{{-}}\NormalTok{EQUIVALENCE(p₁, p₂) then}
            \ControlFlowTok{return}\NormalTok{ false}
\NormalTok{        end }\ControlFlowTok{if}
\NormalTok{    end }\ControlFlowTok{for}
    
    \ControlFlowTok{return}\NormalTok{ true}
\NormalTok{end function}
\end{Highlighting}
\end{Shaded}

\subsection{Specified Performance
Monitor}\label{specified-performance-monitor}

\begin{Shaded}
\begin{Highlighting}[]
\KeywordTok{class}\NormalTok{ PerformanceMonitor:}
    \KeywordTok{def}\NormalTok{ monitor\_execution(}\VariableTok{self}\NormalTok{, }
\NormalTok{                         instruction: Instruction, }
\NormalTok{                         context: Context) }\OperatorTok{{-}\textgreater{}}\NormalTok{ Metrics:}
\NormalTok{        start }\OperatorTok{=}\NormalTok{ time.perf\_counter\_ns()}
\NormalTok{        result }\OperatorTok{=} \VariableTok{self}\NormalTok{.execute(instruction, context)}
\NormalTok{        end }\OperatorTok{=}\NormalTok{ time.perf\_counter\_ns()}
        
        \ControlFlowTok{return}\NormalTok{ Metrics(}
\NormalTok{            latency}\OperatorTok{=}\NormalTok{(end }\OperatorTok{{-}}\NormalTok{ start),}
\NormalTok{            resource\_usage}\OperatorTok{=}\VariableTok{self}\NormalTok{.measure\_resources()}
\NormalTok{        )}
\end{Highlighting}
\end{Shaded}

\subsection{Shipped Peephole Optimizer (Empty Pattern
Table)}\label{shipped-peephole-optimizer-empty-pattern-table}

\begin{Shaded}
\begin{Highlighting}[]
\KeywordTok{class}\NormalTok{ PeepholeOptimizer:}
    \KeywordTok{def}\NormalTok{ optimize(}\VariableTok{self}\NormalTok{, instructions: List[Instruction]) }\OperatorTok{{-}\textgreater{}}\NormalTok{ List[Instruction]:}
\NormalTok{        patterns }\OperatorTok{=} \VariableTok{self}\NormalTok{.load\_patterns()}
\NormalTok{        optimized }\OperatorTok{=}\NormalTok{ []}
\NormalTok{        i }\OperatorTok{=} \DecValTok{0}
        \ControlFlowTok{while}\NormalTok{ i }\OperatorTok{\textless{}} \BuiltInTok{len}\NormalTok{(instructions):}
\NormalTok{            matched }\OperatorTok{=} \VariableTok{False}
            \ControlFlowTok{for}\NormalTok{ pattern }\KeywordTok{in}\NormalTok{ patterns:}
                \ControlFlowTok{if}\NormalTok{ pattern.matches(instructions[i:i}\OperatorTok{+}\NormalTok{pattern.length]):}
\NormalTok{                    optimized.extend(pattern.replace(instructions[i:i}\OperatorTok{+}\NormalTok{pattern.length]))}
\NormalTok{                    i }\OperatorTok{+=}\NormalTok{ pattern.length}
\NormalTok{                    matched }\OperatorTok{=} \VariableTok{True}
                    \ControlFlowTok{break}
            \ControlFlowTok{if} \KeywordTok{not}\NormalTok{ matched:}
\NormalTok{                optimized.append(instructions[i])}
\NormalTok{                i }\OperatorTok{+=} \DecValTok{1}
        \ControlFlowTok{return}\NormalTok{ optimized}
\end{Highlighting}
\end{Shaded}

\subsection{Specified Context-Aware Optimization Sketch
(OPTIMIZE-CONTEXT)}\label{specified-context-aware-optimization-sketch-optimize-context}

Programmatic sketch from the main paper's Section 6.3.2, item 2 (future
work; the performance model and the variant generator are future work,
not implemented machinery):

\begin{verbatim}
function OPTIMIZE-CONTEXT(I, C)
    Let M be a context-performance model, R the rewrite rules
    for each context c in C do
        if PREDICT-PERFORMANCE(M, I, c) < THRESHOLD then
            candidates ← GENERATE-VARIANTS(I, R)
            I ← SELECT-BEST(candidates, c)
            UPDATE-MODEL(M, I, c)
        end if
    end for
    return I
end function
\end{verbatim}

\subsection{Shipped Coordination Flow and Fusion Pass
Listings}\label{shipped-coordination-flow-and-fusion-pass-listings}

The coordination flow implemented by \texttt{AgentCoordinator}
(authenticate, authorize, select by capability, plan stages with
fallback, execute, record history), referenced by the main paper's
Section 5.6:

\begin{verbatim}
Input: Agent set A, Task list τs
Output: Coordinated solutions S

function COORDINATE-AGENTS(A, τs)
    for each task T in τs do
        AUTHENTICATE(task issuer)              // SecurityMonitor
        AUTHORIZE(task, role policy)           // SecurityMonitor
        ag ← SELECT-BY-CAPABILITY(A, T)        // CapabilityProfile (Section 3.5)
        plan ← BUILD-EXECUTION-PLAN(T)         // stages + coordination rules
        task-outcome ← success
        for each stage in plan do
            stage-outcome ← EXECUTE-STAGE(stage, ag)
            if stage-outcome failed and fallback(stage) exists then
                stage-outcome ← EXECUTE-FALLBACK(stage)
            end if
            if stage-outcome failed then
                task-outcome ← failed
                break
            end if
        end for
        RECORD-HISTORY(T, task-outcome)
    end for
end function
\end{verbatim}

The shipped instruction-combining (fusion) pass of the main paper's
Section 4.1.3, also in the source module
\texttt{agentflow.core.isa.instruction\_optimization}:

\begin{Shaded}
\begin{Highlighting}[]
\KeywordTok{def}\NormalTok{ combine\_instructions(}\VariableTok{self}\NormalTok{, insts: List[Instruction]) }\OperatorTok{{-}\textgreater{}}\NormalTok{ List[Instruction]:}
\NormalTok{    combined }\OperatorTok{=}\NormalTok{ []}
\NormalTok{    current\_group }\OperatorTok{=}\NormalTok{ []}
    
    \ControlFlowTok{for}\NormalTok{ inst }\KeywordTok{in}\NormalTok{ insts:}
        \ControlFlowTok{if} \VariableTok{self}\NormalTok{.can\_combine(current\_group, inst):}
\NormalTok{            current\_group.append(inst)}
        \ControlFlowTok{else}\NormalTok{:}
            \ControlFlowTok{if}\NormalTok{ current\_group:}
\NormalTok{                combined.append(}\VariableTok{self}\NormalTok{.merge\_group(current\_group))}
\NormalTok{            current\_group }\OperatorTok{=}\NormalTok{ [inst]}
            
    \ControlFlowTok{if}\NormalTok{ current\_group:}
\NormalTok{        combined.append(}\VariableTok{self}\NormalTok{.merge\_group(current\_group))}
    \ControlFlowTok{return} \VariableTok{self}\NormalTok{.verify\_combinations(combined)}
\end{Highlighting}
\end{Shaded}

\section{Formal Sketches for Future
Directions}\label{formal-sketches-for-future-directions}

Sketches referenced by the main paper's Section 6.3; all are future work
with no current implementation.

\subsection{Stated Sketches}\label{stated-sketches}

\textbf{Compositional reasoning} (Section 6.3.1, item 2): a
parallel-composition rule

\[
\frac{\{P_1\}S_1\{Q_1\} \quad \{P_2\}S_2\{Q_2\} \quad \text{disjoint}(S_1, S_2)}{\{P_1 \land P_2\}S_1 \parallel S_2\{Q_1 \land Q_2\}}
\]

whose disjointness side condition mirrors the \(\parallel\)-typing rule
of the main paper's Definition 3.6.

\textbf{Capability type theory} (Section 6.3.1, item 3):

\begin{verbatim}
Type Capability = {
  Basic: {Perception, Action, Memory}
  Composite: Basic → Basic
  Emergent: P(Composite)
}
\end{verbatim}

with typing rules

\[
\frac{\Gamma \vdash a : \text{cap}_1 \quad \Gamma \vdash b : \text{cap}_2}
     {\Gamma \vdash a \circ b : \text{cap}_1 \otimes \text{cap}_2}
\]

\textbf{Resource consumption model} (Section 6.3.3, item 2):
\(R(i) = \alpha C_{cpu}(i) + \beta C_{mem}(i) + \gamma C_{io}(i)\) with
platform-specific coefficients.

\textbf{Information flow control} (Section 6.3.6, item 1):

\[
\text{Secure}(P) \iff \forall s_1, s_2: L(s_1) = L(s_2) \implies L(\llbracket P \rrbracket s_1) = L(\llbracket P \rrbracket s_2)
\]

where \(L\) assigns security levels to states and outputs (classical
noninterference, terminating runs only).

\section{Artifact Regeneration
Commands}\label{artifact-regeneration-commands}

Exact command lines mirrored from the repository's
\texttt{REPRODUCING.md}, referenced by the main paper's Artifact
Availability section.

\subsection{Section 5.2--5.5
microbenchmarks}\label{section-5.25.5-microbenchmarks}

\begin{verbatim}
./venv/bin/python examples/benchmarks/isa_microbench.py \
    --experiments cache compiler dynamic ablation --seeds 5
\end{verbatim}

\subsection{Section 5.7 real-gateway cache
run}\label{section-5.7-real-gateway-cache-run}

\begin{verbatim}
./venv/bin/python examples/benchmarks/isa_real_llm_bench.py \
    --tasks 200 --seeds 1 --zipf 0.5 1.0 1.5 --max-tokens 64 \
    --output examples/benchmarks/isa_real_llm_results.json
\end{verbatim}

(requires a reachable gateway; the base URL, API key, and CA certificate
are passed per deployment). Any Anthropic-protocol gateway
works---configure \texttt{ANTHROPIC\_BASE\_URL} and
\texttt{ANTHROPIC\_AUTH\_TOKEN}; the paper's run used a local alias for
a hosted GLM reasoning model:

\begin{verbatim}
./venv/bin/python examples/benchmarks/isa_real_llm_bench.py \
    --tasks 200 --seeds 1 --zipf 0.5 1.0 1.5 --max-tokens 64
\end{verbatim}

\subsection{Section 5.7.1 real-gateway compiler/hot-path
harness}\label{section-5.7.1-real-gateway-compilerhot-path-harness}

\begin{verbatim}
./venv/bin/python examples/paper_real_llm_compiler_test.py \
    --backend real --scale paper
\end{verbatim}

(the \texttt{mock} backend reproduces the offline variant;
\texttt{-\/-scale\ paper} selects the 50-realization, 40-instance
protocol of the paper run).

\subsection{Section 5.8 four-task
demonstration}\label{section-5.8-four-task-demonstration}

\begin{verbatim}
./venv/bin/python examples/system_task_demo.py \
    --seed 42 --files 6 --lines-per-file 400 --repeats 12 \
    --hot-threshold 5 --output examples/system_task_results.json
\end{verbatim}

with test suite
\texttt{pytest\ tests/system/test\_system\_task\_demo.py}.

\subsection{Artifact Contents
Inventory}\label{artifact-contents-inventory}

Per-file inventory of the artifact (mirrored from
\texttt{REPRODUCING.md}):

\begin{itemize}
\tightlist
\item
  the AgentISA implementation (\texttt{agentflow.core.isa});
\item
  the microbenchmark suite
  (\texttt{examples/benchmarks/isa\_microbench.py});
\item
  the real-LLM benchmark
  (\texttt{examples/benchmarks/isa\_real\_llm\_bench.py});
\item
  the system-level task demonstration
  (\texttt{examples/system\_task\_demo.py}, raw results in
  \texttt{examples/system\_task\_results.json});
\item
  the four-task workload characterization harness
  (\texttt{examples/system\_task\_workload.py}, tests in
  \texttt{tests/system/test\_system\_task\_workload.py});
\item
  the real-gateway compiler/hot-path harness
  (\texttt{examples/paper\_real\_llm\_compiler\_test.py}), its
  offline-mock and real-gateway results
  (\texttt{examples/paper\_real\_llm\_compiler\_results\_mock.json},
  \texttt{examples/paper\_real\_llm\_compiler\_results\_real.json}), and
  tests (\texttt{tests/system/test\_real\_llm\_compiler\_test.py});
\item
  the raw results
  (\texttt{examples/benchmarks/isa\_microbench\_results.json} and
  \texttt{examples/benchmarks/isa\_real\_llm\_results.json});
\item
  the 43 multi-agent functional tests
  (\texttt{tests/core/isa/test\_coordination.py},
  \texttt{tests/core/instructions/test\_agent\_coordination.py}).
\end{itemize}

\section{Priced-Comparison Protocol and
Table}\label{priced-comparison-protocol-and-table}

Full protocol of the priced comparison summarized by the main paper's
Section 5.7.1; all costs are MODELED from a pricing overlay on measured
tokens.

\textbf{A priced comparison (costs MODELED---computed from a pricing
overlay rather than observed---on measured tokens).} The harness prices
the run's measured token counts under three modes: \emph{no
optimization} (every realization re-queries at full price),
\emph{AgentISA result caching} (first occurrences only---the measured
behavior of both arms), and a \emph{provider prompt cache} (MODELED:
every realization queries, but repeat prompts bill input tokens as
prompt-cache reads at 0.1x the input price---configurable---mirroring
Anthropic's published prompt-caching read discount; completions always
bill at full price). Two simplifications---omitting Anthropic's 1.25x
cache-write premium and 1024-token minimum cacheable prefix---both favor
the provider-cache baseline, so the comparison is conservative for us.
List prices are parameterized stand-ins (\$3/M input, \$15/M output by
default) because the gateway's actual pricing is unknown. Token counts
are measurements, the pricing overlay is a model, and no dollar-savings
claim is made:

\begin{longtable}[]{@{}
  >{\raggedright\arraybackslash}p{(\linewidth - 8\tabcolsep) * \real{0.2000}}
  >{\raggedright\arraybackslash}p{(\linewidth - 8\tabcolsep) * \real{0.2000}}
  >{\raggedright\arraybackslash}p{(\linewidth - 8\tabcolsep) * \real{0.2000}}
  >{\raggedright\arraybackslash}p{(\linewidth - 8\tabcolsep) * \real{0.2000}}
  >{\raggedright\arraybackslash}p{(\linewidth - 8\tabcolsep) * \real{0.2000}}@{}}
\caption{Measured calls and tokens with a MODELED pricing overlay
(real-gateway run, s = 1.0, 50 realizations; the price column is a
model, not a quote).}\tabularnewline
\toprule\noalign{}
\begin{minipage}[b]{\linewidth}\raggedright
Pricing mode
\end{minipage} & \begin{minipage}[b]{\linewidth}\raggedright
Calls
\end{minipage} & \begin{minipage}[b]{\linewidth}\raggedright
Input tokens (full / cache-read)
\end{minipage} & \begin{minipage}[b]{\linewidth}\raggedright
Output tokens
\end{minipage} & \begin{minipage}[b]{\linewidth}\raggedright
Modeled cost
\end{minipage} \\
\midrule\noalign{}
\endfirsthead
\toprule\noalign{}
\begin{minipage}[b]{\linewidth}\raggedright
Pricing mode
\end{minipage} & \begin{minipage}[b]{\linewidth}\raggedright
Calls
\end{minipage} & \begin{minipage}[b]{\linewidth}\raggedright
Input tokens (full / cache-read)
\end{minipage} & \begin{minipage}[b]{\linewidth}\raggedright
Output tokens
\end{minipage} & \begin{minipage}[b]{\linewidth}\raggedright
Modeled cost
\end{minipage} \\
\midrule\noalign{}
\endhead
\bottomrule\noalign{}
\endlastfoot
No optimization & 50 & 9,018 / 0 & 25,600 & \$0.4111 \\
AgentISA result caching & 22 & 4,027 / 0 & 11,264 & \$0.1810 \\
Provider prompt cache (MODELED) & 50 & 4,027 / 4,991 & 25,600 &
\$0.3976 \\
\end{longtable}

On this run's token mix the result cache cuts modeled cost by
56.0\%---it removes whole calls, completions included---while the
modeled provider prompt cache cuts only 3.3\%, because the reasoning
model's completion tokens dominate and a discount on input reads cannot
touch them; the ranking of the two caching strategies is a function of
the input/output token mix, not a constant. On the deterministic mock at
paper scale (200 tasks \(\times\) 1 seed \(\times\) 3 Zipf levels,
template-accounted tokens, same modes and prices) the result cache cuts
modeled cost by 80.1\%, 81.6\%, and 85.1\% at s = 0.5, 1.0, 1.5, while
the modeled provider prompt cache cuts 28.9--31.4\%---there the short
completions do not dominate, so the input-side discount matters more;
both columns remain MODELED-on-measured in both cases. GPTCache, an
external cache layer, is explicitly out of scope; no head-to-head
against it is claimed.

\section{Four-Task Characterization --- Full
Detail}\label{four-task-characterization-full-detail}

Task inventory, per-task instruction streams, configuration, single-task
walk-through, and both result tables of the main paper's Section 5.8;
numbers identical to the main paper.

The four tasks:

\begin{itemize}
\tightlist
\item
  the original \textbf{service-log triage} (linear, 7 instructions),
  detailed below;
\item
  \textbf{csv-etl-anomaly} (fan-out, 8): CSV shard ingest, schema guard,
  parse, per-column statistics, row signatures and anomaly triage,
  aggregate, write;
\item
  \textbf{report-merge} (wide fan-in, 7): five-report ingest, header
  guard, dedupe, per-line signatures, cacheable consistency check,
  merge, write;
\item
  \textbf{regression-diff} (two-branch convergence, 8): parallel
  baseline/candidate ingests, guard, normalize, diff, classify hunks,
  aggregate, write.
\end{itemize}

\textbf{The triage task's instruction stream} (seven instructions over
sandbox-resident files): \texttt{ingest\_logs} reads six 400-line
sandbox log files (side-effectful); \texttt{check\_syntax} is a
defensive format guard that never rejects here (a \emph{planted dead
guard}, standing in for real pipelines' belt-and-braces checks);
\texttt{normalize\_entries} parses and normalizes lines;
\texttt{extract\_incidents} lifts them into incident records;
\texttt{classify\_signature} triage-classifies each incident (LLM-type,
cacheable, backed by a deterministic call-counting rule stand-in);
\texttt{aggregate\_triage} summarizes by severity;
\texttt{write\_report} writes the report (side-effectful).

\textbf{Configuration.} Seed 42; 6 files \(\times\) 400 lines; 12
repeated realizations of the task; hot-path threshold 5. Every phase is
deterministic and re-runnable from the repository (Artifact
Availability).

\textbf{What the triage walk-through shows.}

\begin{itemize}
\tightlist
\item
  \textbf{Static compilation with structural validation.} NONE compiles
  to 7 effective instructions / 7 dispatches; AGGRESSIVE to 6 / 6
  (14.3\%): the dead guard fuses with the following pure step rather
  than being dropped. The walk-through \emph{asserts the compiled
  stream's dataflow validity}---every step whose output is consumed
  downstream must survive compilation, the only eliminable step being
  the guard---so the shrinkage cannot come from dropping live producers.
\item
  \textbf{Reuse across repeated realizations.} The first realization
  makes 50 (simulated) model calls; all eleven subsequent realizations
  make zero (result cache hit rate 91.7\% on the classification step);
  all 12 reports are byte-identical.
\item
  \textbf{Closed-loop hot-path adoption.} The stream crosses the
  hot-path threshold at the sixth recording (the counter must strictly
  exceed it and recording follows execution, so adoption lands at
  realization threshold+2); the runtime \emph{adopts} the optimized form
  automatically, repeats 7--12 execute the shrunken 6-instruction stream
  (instruction executions per realization:
  \texttt{{[}7,7,7,7,7,7,6,6,6,6,6,6{]}}), and on a fresh workspace the
  optimized form reproduces the report
  byte-for-byte---semantics-preserving on the real task, not only on
  synthetic sequences.
\item
  \textbf{Full-stack determinism.} A completely fresh stack---new
  instructions, new cache, new workspace, same seed---reproduces the
  report byte-for-byte (SHA-256 equality).
\end{itemize}

\textbf{Configuration.} Seed 42; 6 files \(\times\) 400 lines; 12
repeated realizations of the task; hot-path threshold 5. Every phase is
deterministic and re-runnable from the repository (Artifact
Availability).

\textbf{What the triage walk-through shows.}

\begin{itemize}
\tightlist
\item
  \textbf{Static compilation with structural validation.} NONE compiles
  to 7 effective instructions / 7 dispatches; AGGRESSIVE to 6 / 6
  (14.3\%): the dead guard fuses with the following pure step rather
  than being dropped. The walk-through \emph{asserts the compiled
  stream's dataflow validity}---every step whose output is consumed
  downstream must survive compilation, the only eliminable step being
  the guard---so the shrinkage cannot come from dropping live producers.
\item
  \textbf{Reuse across repeated realizations.} The first realization
  makes 50 (simulated) model calls; all eleven subsequent realizations
  make zero (result cache hit rate 91.7\% on the classification step);
  all 12 reports are byte-identical.
\item
  \textbf{Closed-loop hot-path adoption.} The stream crosses the
  hot-path threshold at the sixth recording (the counter must strictly
  exceed it and recording follows execution, so adoption lands at
  realization threshold+2); the runtime \emph{adopts} the optimized form
  automatically, repeats 7--12 execute the shrunken 6-instruction stream
  (instruction executions per realization:
  \texttt{{[}7,7,7,7,7,7,6,6,6,6,6,6{]}}), and on a fresh workspace the
  optimized form reproduces the report
  byte-for-byte---semantics-preserving on the real task, not only on
  synthetic sequences.
\item
  \textbf{Full-stack determinism.} A completely fresh stack---new
  instructions, new cache, new workspace, same seed---reproduces the
  report byte-for-byte (SHA-256 equality).
\end{itemize}

\textbf{The four-task characterization.} All four tasks run the
identical closed-loop battery
(\texttt{examples/system\_task\_workload.py}, \texttt{characterize()},
12 realizations each, hot-path threshold 5, one planted dead guard and
one cacheable LLM-type step each, the latter backed by a deterministic
rule stand-in counting simulated model calls):

\emph{Structure (compile-time):}

\begin{longtable}[]{@{}
  >{\raggedright\arraybackslash}p{(\linewidth - 8\tabcolsep) * \real{0.2000}}
  >{\raggedright\arraybackslash}p{(\linewidth - 8\tabcolsep) * \real{0.2000}}
  >{\raggedright\arraybackslash}p{(\linewidth - 8\tabcolsep) * \real{0.2000}}
  >{\raggedright\arraybackslash}p{(\linewidth - 8\tabcolsep) * \real{0.2000}}
  >{\raggedright\arraybackslash}p{(\linewidth - 8\tabcolsep) * \real{0.2000}}@{}}
\caption{Compile-time structure of the four tasks (12 realizations each;
shrink is the NONE \(\rightarrow\) AGGRESSIVE effective-instruction
reduction).}\tabularnewline
\toprule\noalign{}
\begin{minipage}[b]{\linewidth}\raggedright
Task (shape)
\end{minipage} & \begin{minipage}[b]{\linewidth}\raggedright
Instrs
\end{minipage} & \begin{minipage}[b]{\linewidth}\raggedright
Effective count NONE \(\rightarrow\) AGGRESSIVE
\end{minipage} & \begin{minipage}[b]{\linewidth}\raggedright
Shrink
\end{minipage} & \begin{minipage}[b]{\linewidth}\raggedright
Dead guard
\end{minipage} \\
\midrule\noalign{}
\endfirsthead
\toprule\noalign{}
\begin{minipage}[b]{\linewidth}\raggedright
Task (shape)
\end{minipage} & \begin{minipage}[b]{\linewidth}\raggedright
Instrs
\end{minipage} & \begin{minipage}[b]{\linewidth}\raggedright
Effective count NONE \(\rightarrow\) AGGRESSIVE
\end{minipage} & \begin{minipage}[b]{\linewidth}\raggedright
Shrink
\end{minipage} & \begin{minipage}[b]{\linewidth}\raggedright
Dead guard
\end{minipage} \\
\midrule\noalign{}
\endhead
\bottomrule\noalign{}
\endlastfoot
log triage (linear, 7) & 7 & 7 \(\rightarrow\) 6 & -14.3\% &
\texttt{check\_syntax} \\
csv-etl-anomaly (fan-out, 8) & 8 & 8 \(\rightarrow\) 6 & -25.0\% &
\texttt{check\_schema} \\
report-merge (wide fan-in, 7) & 7 & 7 \(\rightarrow\) 6 & -14.3\% &
\texttt{check\_headers} \\
regression-diff (two-branch, 8) & 8 & 8 \(\rightarrow\) 7 & -12.5\% &
\texttt{check\_formats} \\
\end{longtable}

\emph{Execution (runtime, 12 realizations each):}

\begin{longtable}[]{@{}
  >{\raggedright\arraybackslash}p{(\linewidth - 8\tabcolsep) * \real{0.2000}}
  >{\raggedright\arraybackslash}p{(\linewidth - 8\tabcolsep) * \real{0.2000}}
  >{\raggedright\arraybackslash}p{(\linewidth - 8\tabcolsep) * \real{0.2000}}
  >{\raggedright\arraybackslash}p{(\linewidth - 8\tabcolsep) * \real{0.2000}}
  >{\raggedright\arraybackslash}p{(\linewidth - 8\tabcolsep) * \real{0.2000}}@{}}
\caption{Runtime behavior of the four tasks (12 realizations each;
classification-type steps are backed by deterministic rule
stand-ins).}\tabularnewline
\toprule\noalign{}
\begin{minipage}[b]{\linewidth}\raggedright
Task
\end{minipage} & \begin{minipage}[b]{\linewidth}\raggedright
Calls 1st / re-runs
\end{minipage} & \begin{minipage}[b]{\linewidth}\raggedright
Cache hit
\end{minipage} & \begin{minipage}[b]{\linewidth}\raggedright
Hot-path adopted
\end{minipage} & \begin{minipage}[b]{\linewidth}\raggedright
Reports byte-identical
\end{minipage} \\
\midrule\noalign{}
\endfirsthead
\toprule\noalign{}
\begin{minipage}[b]{\linewidth}\raggedright
Task
\end{minipage} & \begin{minipage}[b]{\linewidth}\raggedright
Calls 1st / re-runs
\end{minipage} & \begin{minipage}[b]{\linewidth}\raggedright
Cache hit
\end{minipage} & \begin{minipage}[b]{\linewidth}\raggedright
Hot-path adopted
\end{minipage} & \begin{minipage}[b]{\linewidth}\raggedright
Reports byte-identical
\end{minipage} \\
\midrule\noalign{}
\endhead
\bottomrule\noalign{}
\endlastfoot
log triage & 50 / 0 & 91.7\% & repeat 7 & yes (12/12 + fresh-stack
SHA) \\
csv-etl-anomaly & 14 / 0 & 91.7\% & repeat 7 & yes (12/12 + fresh-stack
SHA) \\
report-merge & 16 / 0 & 91.7\% & repeat 7 & yes (12/12 + fresh-stack
SHA) \\
regression-diff & 240 / 0 & 91.7\% & repeat 7 & yes (12/12 + fresh-stack
SHA) \\
\end{longtable}

\end{document}
